\SourcePage{567}{570}
\section{Current density in a magnetic field}

We shall obtain the quantum-mechanical expression for the current density of a
charged particle moving in a magnetic field. We start from the relation\footnote{In
this section $\mathbf j$ denotes the electric-current density: the particle-current
density multiplied by the charge $e$.}
\begin{equation}
 \delta H=-\frac1c\int\mathbf j\,\delta\mathbf A\,dV,
 \tag{115.1}
\end{equation}
which specifies the change in the Hamilton function of charges distributed in
space when the vector potential is varied\footnote{The Lagrange function for a
charge in a magnetic field contains the term $(e/c)\mathbf v\mathbf A$, or, if
the charge is regarded as distributed through space,
$(1/c)\int\mathbf j\mathbf A\,dV$. Its change upon variation of $\mathbf A$ is
therefore
\[
 \delta L=\frac1c\int\mathbf j\,\delta\mathbf A\,dV.
\]
An infinitesimal change of the Hamilton function is the negative of the change
of the Lagrange function (see I, \S~40).}. In quantum mechanics this relation
must be applied to the mean value of the charged particle's Hamiltonian:
\begin{equation}
 \overline H=\int\Psi^*\left[
 \frac1{2m}\left(\hat{\mathbf p}-\frac ec\mathbf A\right)^2
 -\frac\mu s\mathbf H\hat{\mathbf s}\right]\Psi\,dV.
 \tag{115.2}
\end{equation}
Varying this expression, and using
$\delta\mathbf H=\Curl{\delta\mathbf A}$, gives
\begin{align}
 \delta\overline H={}&\int\Psi^*\left[-\frac e{2mc}
 (\hat{\mathbf p}\,\delta\mathbf A+\delta\mathbf A\hat{\mathbf p})
 +\frac{e^2}{mc^2}\mathbf A\,\delta\mathbf A\right]\Psi\,dV\notag\\
 &-\frac\mu s\int\Curl{\delta\mathbf A}\cdot
 \Psi^*\hat{\mathbf s}\Psi\,dV.
 \tag{115.3}
\end{align}
We transform the term containing $\hat{\mathbf p}\,\delta\mathbf A$ by
integration by parts:
\[
 \int\Psi^*\hat{\mathbf p}\,\delta\mathbf A\Psi\,dV
 =-i\hbar\int\Psi^*\nabla(\delta\mathbf A\Psi)\,dV
 =i\hbar\int\delta\mathbf A\Psi\nabla\Psi^*\,dV
\]
(as usual, the integral over the surface at infinity vanishes). We likewise
integrate the last term of (115.3) by parts, employing the familiar identity of
vector analysis
\[
 \mathbf a\Curl{\mathbf b}
 =-\Div{\mathbf a\times\mathbf b}
 +\mathbf b\Curl{\mathbf a}.
\]

\SourcePage{568}{571}
The integral of the term containing $\Div{}$ vanishes, leaving
\[
 \int\Psi^*\hat{\mathbf s}\Psi\Curl{\delta\mathbf A}\,dV
 =\int\delta\mathbf A\Curl{(\Psi^*\hat{\mathbf s}\Psi)}\,dV.
\]
Thus the final result is
\begin{align*}
 \delta\overline H={}&-\frac{ie\hbar}{2mc}\int\delta\mathbf A
 (\Psi\nabla\Psi^*-\Psi^*\nabla\Psi)\,dV\\
 &+\frac{e^2}{mc^2}\int\mathbf A\,\delta\mathbf A\Psi\Psi^*\,dV
 -\frac\mu s\int\delta\mathbf A\Curl{(\Psi^*\hat{\mathbf s}\Psi)}\,dV.
\end{align*}
Comparison with (115.1) now yields the following current-density expression:
\begin{equation}
 \mathbf j=\frac{ie\hbar}{2m}
 [ (\nabla\Psi^*)\Psi-\Psi^*\nabla\Psi]
 -\frac{e^2}{mc}\mathbf A\Psi^*\Psi
 +\frac\mu s c\Curl{(\Psi^*\hat{\mathbf s}\Psi)}.
 \tag{115.4}
\end{equation}
We stress that, despite its explicit dependence on the vector potential, this
expression is, as required, wholly unambiguous. This follows immediately by a
direct calculation: together with the gauge transformation of the vector
potential, one must, in accordance with (111.8), also transform the wave
function as prescribed by (111.9).

It is also readily verified that the current (115.4), jointly with the charge
density $\rho=e|\Psi|^2$, obeys the continuity equation, as it must:
\[
 \frac{\partial\rho}{\partial t}+\Div{\mathbf j}=0.
\]
The final term in (115.4) is the contribution to the current density due to the
particle's magnetic moment. It has the form $c\Curl{\mathbf m}$, where
\begin{equation}
 \mathbf m=\frac\mu s\Psi^*\hat{\mathbf s}\Psi
 =\Psi^*\hat{\boldsymbol\mu}\Psi
 \tag{115.5}
\end{equation}
is the spatial density of magnetic moment.

Expression (115.4) is the mean value of the current density. It may be viewed
as the diagonal matrix element of an operator, namely the current-density
operator $\hat{\mathbf j}$. This operator is written most simply in the
second-quantized representation: one replaces $\Psi$ and $\Psi^*$ by the
operators $\hat\Psi$ and $\hat\Psi^+$ (and, by the general rule, places
$\hat\Psi^+$ to the left of $\hat\Psi$ in every term). The off-diagonal matrix
elements of this operator may also be defined:
\begin{equation}
 \mathbf j_{nm}=\frac{ie\hbar}{2m}
 [ (\nabla\Psi_n^*)\Psi_m-\Psi_n^*\nabla\Psi_m]
 -\frac{e^2}{mc}\mathbf A\Psi_n^*\Psi_m
 +\frac\mu s c\Curl{(\Psi_n^*\hat{\mathbf s}\Psi_m)}.
 \tag{115.6}
\end{equation}
